\documentclass[11pt]{amsart}

\usepackage[T1]{fontenc}
\usepackage{lmodern}
\usepackage{microtype}
\setlength{\emergencystretch}{2em}
\usepackage[margin=1.15in]{geometry}
\usepackage{amsmath,amssymb,mathtools}
\usepackage{mathrsfs}
\usepackage{amsthm}
\usepackage{booktabs}
\usepackage{enumitem}
\usepackage{xcolor}
\usepackage[colorlinks=true,linkcolor=blue!55!black,urlcolor=blue!55!black]{hyperref}

\newtheorem{theorem}{Theorem}[section]
\newtheorem{proposition}[theorem]{Proposition}
\newtheorem{lemma}[theorem]{Lemma}
\newtheorem{corollary}[theorem]{Corollary}
\newtheorem{definition}[theorem]{Definition}
\newtheorem{example}[theorem]{Example}
\newtheorem{remark}[theorem]{Remark}

\newcommand{\C}{\mathbb C}
\newcommand{\End}{\operatorname{End}}
\newcommand{\Hom}{\operatorname{Hom}}
\newcommand{\rank}{\operatorname{rank}}
\newcommand{\id}{\operatorname{id}}

\title[Resonance-Marked Naturality]{Resonance-Marked Naturality and Homotopy-Tilt Non-Invariance\\
in the Inverse-Leibniz Cubic Case}
\author{Akari H.}
\date{Paper II, manuscript closeout v0.08 --- September 2026}
\hypersetup{citecolor=blue!55!black,pdftitle={Resonance-Marked Naturality and Homotopy-Tilt Non-Invariance},
pdfauthor={Akari H.},pdfsubject={Paper II: normalized tilts and marked boundary responses}}
\setcounter{tocdepth}{1}

\begin{document}

\begin{abstract}
We determine which features of a fixed cubic inverse-Leibniz calculation survive changes
of its normalized homotopy splitting. In a seven-dimensional recurrent envelope,
the full internal tilt orbit consists of all operators satisfying $A(E)=-2E$.
Its common characteristic and minimal polynomial divisor is $S+2$; the complementary
quintic factor of the chosen contraction can be removed. We then define a primitive-free
resonant boundary-response line and a response covector on the tangent subspace for
which the bracket descends through the primitive quotient. We characterize this domain as the kernel of an explicit descent obstruction;
any output quotient extending a transverse cubic direction must erase the marked line. The response is transported by strict comparisons preserving the marked
line, differential, and bracket. An explicit four-generator DGLA has a contractible
resonant pair that a strict quasi-isomorphism cancels, while the Maurer--Cartan functor
over local Artin algebras remains a smooth two-parameter germ. Thus the positive
comparison theorem requires the resonance marking, which the ordinary deformation
functor does not retain. All spectral claims concern the stated internal orbit;
the finite cubic input and the available computational checks are identified explicitly.
\end{abstract}

\maketitle
% The short companion paper keeps its section headings but needs no contents page.

\section{Scope and main conclusions}

The inverse-Leibniz cubic calculation produces two kinds of output.  The first kind belongs to a chosen contraction: a controller box, a recurrent matrix, a cyclic source generator, and a polynomial such as
\[
m_A(S)=(S+2)q(S).
\]
The second kind is forced by the original cochain identities together with a normalized boundary direction.  The purpose of this paper is to separate these two layers and to identify the correct comparison statement.

The main conclusions are:

\begin{enumerate}[label=(\roman*)]
\item the exact recurrent operator is not invariant under admissible homotopy tilts;
\item the normalized \(E\)-rail and its \(-2\) response are invariant inside the fixed envelope;
\item the primitive-free resonant boundary-response pair is natural under strict resonance-faithful comparisons;
\item ordinary DGLA quasi-isomorphism does not preserve that pair, even when the formal Maurer--Cartan deformation behavior is unchanged.
\end{enumerate}

The paper therefore ends with the boxed scope statement
\[
\boxed{\text{strict/resonance-faithful naturality}\;+
       \text{ordinary quasi-isomorphism no-go}.}
\]

No claim is made here for arbitrary triangular systems with spectral multiplicity vector
\[
m_A(S)=\prod_\lambda(S-\lambda)^{e_\lambda}.
\]

\subsection*{Relation to existing deformation theory}
The fixed action-reconstruction problem and its selected contraction are developed
in Paper~I~\cite{PaperI}. The deformation functors used for comparison are standard
Maurer--Cartan constructions \cite{Manetti2005,Getzler2009}. Homotopy transfer changes
representatives and higher operations \cite[Chapter~10]{LodayVallette2012}; it does
not by itself assert invariance of the full matrix $-hD_\alpha$.
Here the additional information is a specified boundary line and its response.
The explicit cancellation example below proves the failure of forgetting this mark
without appealing to a general equivalence theorem.

Paper~III~\cite{PaperIII} extends the polynomial-floor calculation to finite marked
presentations. The present cubic result supplies its motivating example. The two complete
contractions in the v0.08 Paper~I comparison agree on the saved low-rail data;
they illustrate finite higher-transfer dependence, not an identification of
those two completions with every parameter in the internal tilt orbit.

\section{The fixed cubic recurrent envelope}\label{sec:fixed-input}
Throughout this paper, $\alpha,\xi,E$ denote the Hochschild cochains
$\mathsf a,\mathsf x,\mathsf e$ of Paper~I, not the axis perturbations
before their Hochschild lift. The symbol $E$ here is a cochain, not the
polynomial $S+2$. The differential and cup-commutator bracket are those
of Paper~I's constrained Hochschild DGLA.
Its subsection \emph{The presentation quartic class does not lift intrinsically}
and its low-order source appendix specify $\mathsf k,\mathsf p$.
For later use, the source dictionary is
\[
\begin{gathered}
g_3(k)=x_3(k+1)+2x_3(k),\qquad g_3(0)=2\mathsf k-\mathsf p,\\
Y_{13}=[\alpha,\mathsf k]-2[\xi,E],\qquad
C_{13}=Y_{13}-d\mathsf p .
\end{gathered}
\]
Thus $C_{13}$ is a nonclosed complement vector, while $dE$ is a
nonzero boundary. The ordered cochain data and exact replay entry points
are identified in the companion's \path{EVIDENCE_MAP.md}.

Let
\[
C^1\xrightarrow{d}C^2\xrightarrow{d}C^3
\]
be the relevant constrained Hochschild complex, and let
\[
D_\alpha=[\alpha,-]:C^1\longrightarrow C^2
\]
be the operator determined by the cocycle \(\alpha\) tangent to the reduced branch.  A selected partial contraction supplies a finite-dimensional recurrent envelope
\[
\mathcal R\subset C^1,
\qquad
A=-hD_\alpha\big|_{\mathcal R}:\mathcal R\longrightarrow\mathcal R.
\]

More explicitly, for this selected homotopy the ordered recurrent basis is
\[
(E,\ g_3(0),\ Ag_3(0),\ldots,A^5g_3(0)).
\]
These seven cochains are independent and span the invariant envelope
$\mathcal R=\C E\oplus\C[A]g_3(0)$. The six-dimensional cyclic summand
and its companion matrix are established in Paper~I's section
\emph{Triangular spectral extension and residual modules}.
The exact cubic records \cite{CubicRecords} verify the extra $E$ direction
and the source intersections, giving
\[
\dim\mathcal R=7,
\qquad
\rank(D_\alpha|_{\mathcal R})=7,
\]
and
\[
D_\alpha(\mathcal R)\cap Z^2
=D_\alpha(\mathcal R)\cap B^2
=\C D_\alpha E.
\]
The normalized direction satisfies
\[
D_\alpha E=2dE=-2\,\xi\smile\xi,
\qquad
A(E)=-2E.
\]

For the fixed contraction, the characteristic polynomial is
\[
\chi_A(S)=(S+2)^2q(S),
\qquad
q(S)=S^5-4S^4+12S^3-32S^2+80S-192.
\]
This formula is exact, but its status is deliberately local: it describes the selected partial contraction before quotienting by homotopy tilts.


\begin{proposition}[Boundary-compatible completion of the finite source data]\label{prop:completion}
Let \(U=U_{16}\), let \(K_{\mathrm{old}}\subset C^2\) be the fixed
\(100\)-dimensional provenance source space, let \(F=D_\alpha(U)\), and put
\(L=K_{\mathrm{old}}+F\). The reconstructed exact data satisfy
\[
\begin{gathered}
\dim U=\rank(d|_U)=35,\quad \dim F=23,\quad \dim L=116,\\
\dim(F\cap K_{\mathrm{old}})=7,\qquad
L\cap B^2=K_{\mathrm{old}}\cap B^2=\C dE.
\end{gathered}
\]
The fixed map \(h_{\mathrm{old}}:K_{\mathrm{old}}\to U\) has \(h_{\mathrm{old}}(dE)=E\).
Its extensions to \(L\), with values in \(U\), form an affine space modeled on
\[
\Hom(L/K_{\mathrm{old}},U),
\]
of dimension \(16\cdot35=560\). Each has a boundary-compatible extension to a
standard contraction, allowing the remaining splitting choices to vary.
\end{proposition}
\begin{proof}
The ranks and intersections are recomputed from the cochains by the v0.03
supplement; the full constrained boundary space has dimension \(88\).
A linear extension is free on a complement of \(K_{\mathrm{old}}\) in \(L\).
Since \(d|_U\) is injective and \(L\cap dU=\C dE\), every such extension
agrees on that intersection with \((d|_U)^{-1}:dU\to U\).
It therefore extends consistently to \(L+dU\) with \(hd|_U=\id\).
Choose the degree-one acyclic complement to contain \(U\), extend the inverse
on boundaries, and choose cycle representatives inside the resulting kernel of
\(h\). On a closed element \(z\in L\), the vector \(z-dhz\) is closed; the
one-dimensional boundary intersection ensures the required cycle representatives
are independent modulo boundaries. Extending these vector-space splittings
degree by degree gives a standard contraction.
No previously unspecified cohomology representatives are held fixed in this assertion.
\end{proof}

\begin{remark}
Cycles must not be confused with boundaries here: the recomputed values are
\(\dim(L\cap Z^2)=2\) and \(\dim(L\cap B^2)=1\).
The additional cohomology direction is compatible with the splitting construction;
it cannot be discarded when checking contraction identities.
\end{remark}

\section{The exact homotopy-tilt action}

\begin{definition}[Normalized source directions]
Put
\[
F=D_\alpha(\mathcal R),
\qquad
F_{\mathrm{norm}}=F\cap Z^2,
\qquad
\overline F=F/F_{\mathrm{norm}}.
\]
The normalized internal tilt group is
\[
\mathbb T_{\mathcal R}=\Hom(\overline F,\mathcal R).
\]
\end{definition}

For \(\tau\in\mathbb T_{\mathcal R}\), write \(q_F:F\to\overline F\) for the quotient map and define
\[
h_\tau(q)=h(q)+\tau(q_F(q)),
\qquad q\in F.
\]
The normalized directions are fixed because \(q_F\) vanishes on \(F_{\mathrm{norm}}\).  The residual source changes by
\[
r_\tau(q)=q-dh_\tau(q),
\qquad
dr_\tau(q)=dq.
\]
Thus the degree-three differential data of the nonclosed source quotient are unchanged.

\begin{proposition}[Exact operator action]
The recurrent operator changes according to
\[
A_\tau=A-\tau q_FD_\alpha|_{\mathcal R}.
\]
The map
\[
\mathbb T_{\mathcal R}
\longrightarrow
\{T\in\End(\mathcal R):T(E)=0\},
\qquad
\tau\longmapsto \tau q_FD_\alpha|_{\mathcal R},
\]
is an isomorphism in the cubic model.
\end{proposition}

\begin{proof}
The first identity is immediate from \(A=-hD_\alpha\) and the definition of \(h_\tau\).  Since \(D_\alpha|_{\mathcal R}\) has rank seven while \(F_{\mathrm{norm}}=\C D_\alpha E\), the quotient map \(q_FD_\alpha\) has kernel \(\C E\) and identifies
\[
\mathcal R/\C E\cong\overline F.
\]
Every endomorphism vanishing on \(E\) therefore factors uniquely through \(q_FD_\alpha\).
\end{proof}

\begin{lemma}[The internal tilts are realizable splittings]
Every tilt in the displayed internal family extends to a contraction of the
full constrained complex. Only its stated normalization is fixed; the family
does not require retention of all $100$ archived homotopy assignments.
\end{lemma}
\begin{proof}
The recurrent envelope lies in the $35$-dimensional space $U$ on which $d$
is injective. Choose a degree-one acyclic complement containing $U$.
The variation $\tau q_F$ on $F$ vanishes on $F\cap Z^2$ and has image in
that complement. Extend it by zero on $Z^2$ and then linearly to $C^2$.
Adding it to a completed homotopy preserves its inverse on boundaries and
its zero values on the chosen cohomology representatives.
The resulting map $h'$ has image in the same acyclic complement and is
surjective onto it. Inside $\ker h'$, extend the chosen cohomology
representatives by a complement $W$; because
$\ker h'\cap Z^2$ is exactly those representatives, $W$ is a complement to
$Z^2$. Projection along $B^2\oplus W$ and the adjacent-degree splittings
give the required contraction.
\end{proof}

\begin{corollary}[Cubic tilt orbit]
In the fixed seven-dimensional envelope,
\[
\dim\mathbb T_{\mathcal R}=6\cdot7=42,
\]
and
\[
\boxed{
\mathcal O_{\mathrm{tilt}}(A)
=\{B\in\End(\mathcal R):B(E)=-2E\}.}
\]
\end{corollary}

\begin{proof}
The dimension follows from \(\dim\overline F=6\) and \(\dim\mathcal R=7\).  The preceding proposition says precisely that all variations vanishing on \(E\) occur, while the normalized relation \(A(E)=-2E\) is fixed.
\end{proof}

The number $42$ here is the dimension of a parameter space
$\Hom(\mathcal R/\C E,\mathcal R)$. It is not the full cohomology space
$H^2$ of Paper~I, whose separately computed dimension happens also to be $42$.
Nor is this $42$-parameter family the $560$-parameter extension family fixing
all of $K_{\mathrm{old}}$ in Proposition~\ref{prop:completion}.

\section{Spectral non-invariance and the surviving rail}

The tilt orbit contains the two exact representatives
\[
B_0=\operatorname{diag}(-2,0,0,0,0,0,0),
\qquad
B_3=\operatorname{diag}(-2,3,3,3,3,3,3).
\]
Their characteristic polynomials are
\[
\chi_{B_0}(S)=(S+2)S^6,
\qquad
\chi_{B_3}(S)=(S+2)(S-3)^6.
\]
Both operators agree with \(A\) on \(\C E\) and hence are exact normalized internal tilts.  We obtain:

\begin{theorem}[Homotopy-tilt spectral no-go]
The factor \(q(S)\), the second \(-2\) eigenline, the six-dimensional cyclic module generated by the fixed \(g_3(0)\), and the fixed-contraction recurrence \((S+2)q(S)\) are not invariants of the original cubic deformation problem under normalized internal homotopy tilts.  The greatest common divisor of the characteristic polynomials over the full internal tilt orbit is
\[
\gcd_{\tau}\chi_{A_\tau}(S)=S+2,
\]
and the same statement holds for minimal polynomials.
\end{theorem}

\begin{proof}
The two displayed representatives already change the complementary spectrum.  Conversely every operator in the orbit has the eigenrelation \(B(E)=-2E\), so \(S+2\) divides every characteristic and minimal polynomial.  Since the induced operator on \(\mathcal R/\C E\) is arbitrary, no nonconstant factor on that quotient can be common to the whole orbit.
\end{proof}

The surviving temporal response is therefore the normalized \(E\)-rail:
\[
\boxed{
D_\alpha A^nE=2(-2)^n dE
\qquad(n\geq0).}
\]
This is a chain-level boundary response, not the full spectrum of a chosen recurrent matrix.

\begin{remark}
The dimension seven in this section is an envelope dimension for the present internal tilt calculation.  It is not asserted to be invariant under an arbitrary change of partial contraction that changes the recurrent envelope itself.
\end{remark}

\section{Dependence of the cyclic source on the splitting}

The fixed low-rail construction has
\[
g_3(0)=2\mathsf k-\mathsf p.
\]
At the first nontrivial source stage, the admissible triangular splitting tilt
\[
\mathsf p_t=\mathsf p+t w,
\qquad
C_{13,t}=C_{13}-tdw,
\]
leaves the original source equation and its differential-rank placement unchanged.  Taking \(w=E\) gives
\[
g_{3,t}(0)=g_3(0)-tE,
\]
and hence
\[
D_\alpha g_{3,t}(0)
=D_\alpha g_3(0)-tD_\alpha E.
\]
Since \(D_\alpha E\neq0\), the exact \(g_3\) source sequence is not determined by the original equations independently of the contraction.  Taking \(w=\mathsf k\) gives the stronger variation
\[
D_\alpha g_{3,t}(0)
=D_\alpha g_3(0)-t[\alpha,\mathsf k],
\qquad
d[\alpha,\mathsf k]\neq0,
\]
so this ambiguity is not merely a boundary ambiguity.

In contrast, the \(E\)-rail is forced by the original identities
\[
dE=-\xi\smile\xi,
\qquad
[\alpha,E]=-2\xi\smile\xi=2dE,
\]
and the normalized contraction condition \(h(dE)=E\).  Therefore
\[
A(E)=-2E
\]
is stable under the normalized internal action.

\section{The primitive-free resonant response object}

The \(E\)-dependent description is useful for computation but not optimal for comparison.  We now remove the primitive choice.

\begin{definition}[Resonant boundary-response pair]
Let
\[
\beta=\xi\smile\xi,
\qquad
\mathcal L_{\mathrm{res}}=\C\beta\subset B^2,
\]
and define
\[
\mathcal P_{\mathrm{res}}
=d^{-1}(\mathcal L_{\mathrm{res}})/Z^1.
\]
The differential induces an isomorphism
\[
\bar d:\mathcal P_{\mathrm{res}}\xrightarrow{\sim}\mathcal L_{\mathrm{res}}.
\]
\end{definition}

Let $T=H^1(C^\bullet)$ with specified cocycle representatives.
Since this constrained complex has $C^0=0$, we identify it with $Z^1$.
In the cubic model $T=\C\alpha\oplus\C\xi$ is the two-dimensional
Zariski tangent space of $\C[[u,v]]/(v^3)$, not the one-dimensional
tangent space of its reduced scheme $\C[[u]]$.
Define the linear subspace of response directions
\[
T_{\mathrm{resp}}=
\{a\in T:
[a,Z^1]=0,\ 
[a,d^{-1}(\mathcal L_{\mathrm{res}})]\subseteq\mathcal L_{\mathrm{res}}\}.
\]
For \(a\in T_{\mathrm{resp}}\), and only under these conditions, \(D_a=[a,-]\) descends to
\[
\bar D_a:\mathcal P_{\mathrm{res}}\longrightarrow\mathcal L_{\mathrm{res}}.
\]

\begin{definition}[Response operator and covector]
For \(a\in T_{\mathrm{resp}}\), define
\[
\Theta(a)=\bar D_a\bar d^{-1}
\in\End(\mathcal L_{\mathrm{res}}).
\]
Since \(\mathcal L_{\mathrm{res}}\) is one-dimensional, there is a unique scalar \(\rho_{\mathrm{res}}(a)\) such that
\[
\Theta(a)=\rho_{\mathrm{res}}(a)\id.
\]
Thus
\[
\rho_{\mathrm{res}}\in T_{\mathrm{resp}}^\vee.
\]
For the fixed cubic complex of Paper~I, \(Z^1=\C\alpha\oplus\C\xi\).
The exact cochains satisfy
\[
[\alpha,\alpha]=[\alpha,\xi]=0,\quad
[\xi,\xi]=2\beta\ne0,\quad [\alpha,E]=2dE.
\]
Since \(d^{-1}(\mathcal L_{\mathrm{res}})=\C E+Z^1\), these identities prove
\(T_{\mathrm{resp}}=\C\alpha\) and \(\rho_{\mathrm{res}}(\alpha)=2\).
Thus the response covector is defined on a proper subspace of the two-dimensional
tangent space. The extra descent condition cannot be suppressed.
\end{definition}


\begin{theorem}[Maximal descent domain and obstruction]\label{thm:response-obstruction}
Put \(Z=Z^1\), \(V_{\mathrm{res}}=d^{-1}(\mathcal L_{\mathrm{res}})\),
and let \(\pi:C^2\to C^2/\mathcal L_{\mathrm{res}}\) be the quotient.
For the specified tangent representatives, the linear obstruction map
\[
\mathfrak o:T\longrightarrow
\Hom(Z,C^2)\oplus\Hom(V_{\mathrm{res}},C^2/\mathcal L_{\mathrm{res}}),
\qquad
a\longmapsto\bigl(D_a|_Z,\pi D_a|_{V_{\mathrm{res}}}\bigr)
\]
satisfies \(T_{\mathrm{resp}}=\ker\mathfrak o\).
This is the largest linear domain on which the unmodified bracket gives a
primitive-independent \(\mathcal L_{\mathrm{res}}\)-valued response.
A strict isomorphism preserving the full tangent representatives and marked
boundary line transports this kernel, as well as its response covector.

In the cubic model, if \(a=s\alpha+t\xi\) with \(t\ne0\), any linear map
\(Q:C^2\to W\) for which \(QD_a\) descends through \(V_{\mathrm{res}}/Z\)
must kill the entire resonant line:
\[
Q(\beta)=0.
\]
Thus quotienting the output cannot extend this primitive-free response to such
a direction while retaining a nonzero image of the same marked line.
\end{theorem}
\begin{proof}
A map \(D_a:V_{\mathrm{res}}\to C^2\) induces a map from \(V_{\mathrm{res}}/Z\)
to \(\mathcal L_{\mathrm{res}}\) exactly when it kills \(Z\) and has image in
\(\mathcal L_{\mathrm{res}}\). These are precisely the two components of
\(\mathfrak o(a)=0\); bilinearity proves linearity in \(a\).
Preservation of differential and bracket transports both conditions.
For the cubic claim,
\[
D_a\xi=[s\alpha+t\xi,\xi]=2t\beta.
\]
Since \(\xi\in Z\), descent requires \(2tQ(\beta)=0\); over \(\C\) this gives
the conclusion. Conversely the verified identities show that every \(s\alpha\)
satisfies both conditions. Therefore \(T_{\mathrm{resp}}=\C\alpha\) is exact,
not merely a sufficient domain.
\end{proof}

\section{Strict resonance-faithful naturality}

\begin{definition}[Resonance-faithful comparison]
Let two constrained DGLA models carry marked data
\[
(T_{\mathrm{resp}},\mathcal L_{\mathrm{res}},\mathcal P_{\mathrm{res}},\bar d,\bar D)
\quad\text{and}\quad
(T'_{\mathrm{resp}},\mathcal L'_{\mathrm{res}},\mathcal P'_{\mathrm{res}},\bar d',\bar D').
\]
A comparison consists of isomorphisms
\[
\Phi_T:T_{\mathrm{resp}}\xrightarrow{\sim}T'_{\mathrm{resp}},
\quad
\Phi_L:\mathcal L_{\mathrm{res}}\xrightarrow{\sim}\mathcal L'_{\mathrm{res}},
\quad
\Phi_P:\mathcal P_{\mathrm{res}}\xrightarrow{\sim}\mathcal P'_{\mathrm{res}},
\]
such that
\[
\bar d'\Phi_P=\Phi_L\bar d,
\qquad
\bar D'_{\Phi_Ta}\Phi_P=\Phi_L\bar D_a.
\]
If the comparison comes from a strict DGLA isomorphism, these identities follow from preservation of the differential and bracket together with preservation of the resonance marking.
\end{definition}

\begin{theorem}[Naturality of the marked response]
Under a resonance-faithful comparison,
\[
\boxed{
\Theta'(\Phi_Ta)=\Phi_L\Theta(a)\Phi_L^{-1}.}
\]
Consequently,
\[
\boxed{
\rho'_{\mathrm{res}}(\Phi_Ta)=\rho_{\mathrm{res}}(a),}
\]
or equivalently
\[
\rho'_{\mathrm{res}}=(\Phi_T^{-1})^\vee\rho_{\mathrm{res}}.
\]
\end{theorem}

\begin{proof}
Using the two commuting identities gives
\[
\begin{aligned}
\Theta'(\Phi_Ta)
&=\bar D'_{\Phi_Ta}(\bar d')^{-1}\\
&=\Phi_L\bar D_a\Phi_P^{-1}
       \bigl(\Phi_P\bar d^{-1}\Phi_L^{-1}\bigr)\\
&=\Phi_L\Theta(a)\Phi_L^{-1}.
\end{aligned}
\]
The resonant line is one-dimensional, so conjugation leaves its scalar unchanged.  The covector formula is the same statement in dual form.
\end{proof}

\begin{remark}
Rescaling the normal generator changes \(\beta\) and hence the chosen basis of \(\mathcal L_{\mathrm{res}}\), but the conjugation formula cancels that scalar change.  The line and the response operator, rather than a fixed boundary vector, are the intrinsic marked objects.
\end{remark}

\section{Ordinary quasi-isomorphism no-go}

\begin{theorem}[Contractible-pair cancellation]
There are strict DGLA quasi-isomorphisms that preserve the formal Maurer--Cartan deformation behavior but erase the marked resonant boundary-response pair.
\end{theorem}

\begin{proof}
Let \(\mathfrak g\) have degree-one generators \(\alpha,\xi,E\) and degree-two generator \(\beta\), with
\[
d\alpha=d\xi=0,
\qquad
dE=\beta,
\]
and the only nonzero brackets
\[
[\xi,\xi]=2\beta,
\qquad
[\alpha,E]=2\beta.
\]
Let \(\mathfrak h\) be the abelian DGLA on degree-one generators \(\alpha,\xi\), and let
\[
p:\mathfrak g\longrightarrow\mathfrak h
\]
kill \(E\) and \(\beta\) while fixing \(\alpha\) and \(\xi\).
All brackets land in the central degree-two line, so the graded Jacobi identity holds.
The differential is a derivation because all resulting degree-three terms vanish.
The kernel is the acyclic ideal \(E\xrightarrow{d}\beta\); hence \(p\) is a strict
quasi-isomorphism. Both models have two-dimensional degree-one cohomology and
vanishing degree-two cohomology.

For \(\gamma=u\alpha+v\xi+wE\), the Maurer--Cartan equation in \(\mathfrak g\) is
\[
w+v^2+2uw=0,
\]
over a local Artin \(\C\)-algebra with \(u,v,w\) in its maximal ideal.
Since \(1+2u\) is a unit, \(w=-v^2/(1+2u)\) is uniquely determined.
The assignment is natural in the Artin algebra. There are no degree-zero gauge
directions, so this directly identifies the two Maurer--Cartan functors with the
same smooth two-parameter germ.

For the response construction,
\(Z^1=\C\alpha\oplus\C\xi\) and \(T_{\mathrm{resp}}=\C\alpha\):
\([\alpha,Z^1]=0\), whereas \([\xi,\xi]=2\beta\ne0\) excludes the \(\xi\)-direction.
Thus the response covector is defined on this one-dimensional subspace.
In \(\mathfrak g\) one has
\[
\mathcal L_{\mathrm{res}}=\C\beta,
\qquad
\rho_{\mathrm{res}}(\alpha)=2,
\]
whereas in \(\mathfrak h\) the product \([\xi,\xi]\) vanishes and
\[
\mathcal L'_{\mathrm{res}}=0.
\]
Thus the ordinary quasi-isomorphism erases precisely the contractible chain-level datum that the marked theory is designed to retain.
\end{proof}

\begin{corollary}
The pair \((\mathcal L_{\mathrm{res}},\rho_{\mathrm{res}})\) is not an invariant of the ordinary DGLA derived category or of the associated formal moduli problem alone.
\end{corollary}

\section{Computational evidence and scope}\label{sec:verification}

The finite cubic input is recorded in \cite{CubicRecords} and is supported by exact
rational linear algebra. The locally rerunnable part of the verification family checks:

\begin{itemize}
\item the support/incidence model and the \(35\)-dimensional \(U_{16}\) controller image;
\item the \(19\)-dimensional forced domain and \(7\)-dimensional recurrent image;

\item the \(42\)-dimensional internal tilt module and the exact orbit formula;
\item the two contrasting characteristic polynomials proving spectral non-invariance.
\end{itemize}

The bundled supplement reconstructs the Hochschild differentials from the finite
algebra action and v0.33 cochain certificate, independently of the unavailable
v0.48 JSON and v0.53 helper. It verifies \(\dim C^1_B=90\), \(\rank d_1=88\),
the two-dimensional cocycle basis, and the source/boundary intersections in
Proposition~\ref{prop:completion}. It also checks the three \(g_3\)-source identities,
the nonzero boundary \(E\)-tilt, \(d[\alpha,\mathsf k]\ne0\), and
\(d^2\mathsf k=0\), together with typed naturality and the cancellation model.
Thus these claims have self-contained replacement checks. The all-\(n\)
\(E\)-rail identity follows from \(A(E)=-2E\), not from finite samples.

The following questions are intentionally excluded from Paper II:

\begin{enumerate}[label=(\arabic*)]
\item a general theorem for arbitrary \(m_A(S)=\prod_\lambda(S-\lambda)^{e_\lambda}\);

\item contraction-independent classification of the full \(U_{16}\) envelope;
\item a canonical gauge-fixing that recovers the richer \(q\)-spectrum;
\item simultaneous deformation of the algebra \(A\) and the tensor \(B\).
\end{enumerate}

The response covector is therefore functorial for the specified resonance-faithful
comparisons on \(T_{\mathrm{resp}}\). The cancellation example proves that its
nonzero boundary line cannot be recovered from an ordinary quasi-isomorphism class
alone. The cubic polynomial floor and this comparison theorem answer different
parts of the same question: what data remain after the permitted changes of
splitting and presentation?

\clearpage
\begin{thebibliography}{99}
\raggedright

\bibitem{Manetti2005}
M.~Manetti,
\emph{Deformation theory via differential graded Lie algebras},
expository notes (written 1999), arXiv:math/0507284, 2005.
\url{https://arxiv.org/abs/math/0507284}.

\bibitem{Getzler2009}
E.~Getzler,
\emph{Lie theory for nilpotent $L_\infty$-algebras},
Ann. of Math. (2) \textbf{170} (2009), 271--301.
\href{https://doi.org/10.4007/annals.2009.170.271}{doi:10.4007/annals.2009.170.271}.

\bibitem{LodayVallette2012}
J.-L.~Loday and B.~Vallette,
\emph{Algebraic Operads}, Grundlehren Math. Wiss. \textbf{346},
Springer, Heidelberg, 2012, Chapter~10.
\href{https://doi.org/10.1007/978-3-642-30362-3}{doi:10.1007/978-3-642-30362-3}.

\bibitem{PaperI}
Akari H.,
\emph{The Inverse Leibniz Problem: Reconstruction Fibers, Rigidity, Obstructions, and Deformation DGLAs},
Paper I, manuscript closeout v0.08, based on the fixed-case v0.41 record, 2026.
Unpublished companion manuscript.

\bibitem{PaperIII}
Akari H.,
\emph{General Spectral Floors from Marked Deformation Presentations},
Paper III, manuscript closeout v0.08, 2026. Unpublished companion manuscript.

\bibitem{CubicRecords}
Akari H.,
\emph{Cubic support, forced-core, source-behavior, and homotopy-tilt records},
versions v0.47--v0.54, 2026. Unpublished computational records.
Replacement checks are included in the companion
\path{three_papers_v0_08_bundle.zip}, under
\path{evidence/legacy/verification/part_II/}, with recovery checks in
\path{evidence/legacy/verify_proof_recovery_v0_03.py}.
See \path{EVIDENCE_MAP.md} for claim-to-file locations and
Section~\ref{sec:verification} for scope. No public archive identifier
has been assigned in this local version.
\end{thebibliography}
\end{document}
